\section*{Implementation}

\subsection*{End-to-end Workflow}
\begin{figure}[H]
\centering
\resizebox{1.08\linewidth}{!}{%
\begin{tikzpicture}[
    node distance=9mm and 10mm,
    box/.style={rectangle, rounded corners, draw=black, thick, align=center, text width=32mm, minimum height=8mm, fill=blue!12},
    readbox/.style={box, fill=blue!12},
    writebox/.style={box, fill=green!18},
    procfs/.style={box, fill=purple!30},
    decision/.style={diamond, draw=black, thick, fill=green!18, aspect=2.0, align=center, inner sep=1.3pt, text width=20mm},
    flow/.style={-Latex, thick, shorten >=1.5pt, shorten <=1.5pt}
]

\node[procfs] (start) {{\scriptsize\texttt{dirty\_tracking\_start}}};
\node[box, below=of start] (downgrade) {Active tracking page table initalized. Cumulative tracking page table initialized. Downgrade mapped GPU pages to \texttt{READ\_ONLY}\\and snapshot pre-tracking permission masks};
\node[box, below=of downgrade] (active) {Tracking active};

\node[procfs, left=36mm of active] (pausecmd) {{\scriptsize\texttt{dirty\_tracking\_pause}}};
\node[box, below=of pausecmd] (paused) {Tracking paused\\(\texttt{active=false})};
\node[procfs, below=of paused] (resume) {{\scriptsize\texttt{dirty\_tracking\_resume}}};

\node[procfs, below=20mm of active] (cutover) {{\scriptsize\texttt{dirty\_tracking\_query\_}}\\[-1pt]{\scriptsize\texttt{cutover}}};
\node[box, below=of cutover] (snapshot) {Snapshot pending\\(live table swapped out)};
\node[procfs, below=of snapshot] (dump) {{\scriptsize\texttt{dirty\_tracking\_query\_}}\\[-1pt]{\scriptsize\texttt{dump}}};

\node[decision, right=70mm of active] (access) {GPU\\fault?};
\node[readbox, below left=10mm and 10mm of access] (readpath) {Read fault path\\Map/retain \texttt{READ\_ONLY}\\No dirty insert};
\node[writebox, below right=10mm and 10mm of access] (writepath) {Write/atomic fault path\\ Map writable permissions \texttt{READ\_WRITE} or \texttt{READ\_WRITE\_ATOMIC}\\Insert page + timestamp};

\node[decision, left=70mm of paused] (pausedaccess) {GPU\\fault?};
\node[readbox, below left=10mm and 10mm of pausedaccess] (pausedreadpath) {Read fault path\\Map/retain \texttt{READ\_ONLY}\\No dirty insert while paused};
\node[writebox, below right=10mm and 10mm of pausedaccess] (pausedwritepath) {Write/atomic fault path\\Map writable permissions\\\texttt{READ\_WRITE} or \texttt{READ\_WRITE\_ATOMIC}\\No dirty insert while paused};

\node[procfs, right=36mm of dump] (stop) {{\scriptsize\texttt{dirty\_tracking\_stop}}};
\node[box, below=of stop] (restore) {Restore saved \texttt{READ\_WRITE}/\texttt{ATOMIC}\\permissions and destroy session state};

\draw[flow] (start) -- (downgrade);
\draw[flow] (downgrade) -- (active);

\draw[flow] (active.west) -- (pausecmd.east);
\draw[flow] (pausecmd) -- (paused);
\draw[flow] (paused) -- (resume);
\draw[flow] (resume.east) to[out=0,in=220] ([xshift=-2mm]downgrade.west);

\draw[flow] (active.south west) to[out=-110,in=110]
    node[pos=0.5, right=1mm, align=center] {\scriptsize \texttt{snapshot\_pending == false}\\\scriptsize{for cutover}}
    (cutover.north west);
\draw[flow] (cutover.north east) to[out=70,in=-70] (active.south east);
\draw[flow] (cutover) -- (snapshot);
\draw[flow] (snapshot) -- (dump);
\draw[flow] (dump.west) to[out=180,in=-110] (cutover.west);

\draw[flow] (active.east) -- (access.west);
\draw[flow] (access.south west) to[out=240,in=50] node[above left, pos=0.45] {\scriptsize read} (readpath.north);
\draw[flow] (access.south east) to[out=300,in=130] node[above right, pos=0.45] {\scriptsize write/atomic} (writepath.north);
\draw[flow] (readpath.north west) to[out=135,in=-45] (active.east);
\draw[flow] (writepath.north east) to[out=45,in=45] (active.north east);

\draw[flow] (paused.west) -- (pausedaccess.east);
\draw[flow] (pausedaccess.south west) to[out=240,in=50] node[above left, pos=0.45] {\scriptsize read} (pausedreadpath.north);
\draw[flow] (pausedaccess.south east) to[out=300,in=130] node[above right, pos=0.45] {\scriptsize write/atomic} (pausedwritepath.north);
\draw[flow] (pausedreadpath.north west) to[out=45,in=150] (paused.west);
\draw[flow] (pausedwritepath.north east) to[out=135,in=-135] (paused.south west);

\draw[flow] (active.east) to[out=0,in=180] (stop.west);
\draw[flow] (stop) -- (restore);

\end{tikzpicture}
}
\caption{End-to-end workflow of the dirty page tracking system, illustrating procfs control paths, GPU fault handling, and state transitions.}
\label{fig:permission-lifecycle-flow}
\end{figure}
\subsubsection*{State Machine}
Dirty tracking is modelled as a single-owner lifecycle state machine, with
all lifecycle operations serialized by
\texttt{uvm\_dirty\_lifecycle\_lock} and guarded by owner-PID checks.
The states are:
\begin{enumerate}
    \item \textbf{Tracking stopped.}
          \texttt{started=false}, \texttt{active=false},
          \texttt{snapshot\_pending=false}.
    \item \textbf{Tracking active with snapshot.}
          \texttt{started=true}, \texttt{active=true},
          \texttt{snapshot\_pending=true}.
    \item \textbf{Tracking active without snapshot.}
          \texttt{started=true}, \texttt{active=true},
          \texttt{snapshot\_pending=false}.
    \item \textbf{Tracking paused with snapshot.}
          \texttt{started=true}, \texttt{active=false},
          \texttt{snapshot\_pending=true}.
    \item \textbf{Tracking paused without snapshot.}
          \texttt{started=true}, \texttt{active=false},
          \texttt{snapshot\_pending=false}.
\end{enumerate}

\begin{figure}[H]
\centering
\resizebox{0.98\linewidth}{!}{%
\begin{tikzpicture}[
    st/.style={rectangle, rounded corners, draw=black, thick, align=center, text width=35mm, minimum height=6mm, fill=blue!10},
    tr/.style={-Latex, thick}
]
\node[st] (stopped) at (0mm,22mm) {\small{Tracking stopped}\\{\scriptsize \texttt{started=false}}};
\node[st] (active_no) at (-46mm,4mm) {\small{Tracking active\\without snapshot}\\{\scriptsize \texttt{started=true, active=true}}\\{\scriptsize \texttt{snapshot\_pending=false}}};
\node[st] (active_yes) at (46mm,4mm) {\small{Tracking active\\with snapshot}\\{\scriptsize \texttt{started=true, active=true}}\\{\scriptsize \texttt{snapshot\_pending=true}}};
\node[st] (paused_no) at (-46mm,-20mm) {\small{Tracking paused\\without snapshot}\\{\scriptsize \texttt{started=true, active=false}}\\{\scriptsize \texttt{snapshot\_pending=false}}};
\node[st] (paused_yes) at (46mm,-20mm) {\small{Tracking paused\\with snapshot}\\{\scriptsize \texttt{started=true, active=false}}\\{\scriptsize \texttt{snapshot\_pending=true}}};

\draw[tr] (stopped.west) to[out=225,in=10] node[left, pos=0.50] {\scriptsize \texttt{start}} (active_no.north);

\draw[tr] (active_no.east) to[bend left=16] node[above] {\scriptsize \texttt{query\_cutover}} (active_yes.west);
\draw[tr] (active_yes.west) to[bend left=16] node[below] {\scriptsize \texttt{query\_dump}} (active_no.east);
\draw[tr] (paused_no.east) to[bend left=16] node[above] {\scriptsize \texttt{query\_cutover}} (paused_yes.west);
\draw[tr] (paused_yes.west) to[bend left=16] node[below] {\scriptsize \texttt{query\_dump}} (paused_no.east);

\draw[tr] (active_no.south) to[bend left=12] node[left] {\scriptsize \texttt{pause}} (paused_no.north);
\draw[tr] (paused_no.north) to[bend left=12] node[right] {\scriptsize \texttt{resume}} (active_no.south);
\draw[tr] (active_yes.south) to[bend right=12] node[right] {\scriptsize \texttt{pause}} (paused_yes.north);
\draw[tr] (paused_yes.north) to[bend right=12] node[left] {\scriptsize \texttt{resume}} (active_yes.south);

\draw[tr] (active_no.north) to[bend left=16] node[left, pos=0.55] {\scriptsize \texttt{stop}} (stopped.west);
\draw[tr] (active_yes.north) to[bend right=16] node[right, pos=0.55] {\scriptsize \texttt{stop}} (stopped.east);
\draw[tr] (paused_no.north west) to[out=150,in=150] node[left, pos=0.38] {\scriptsize \texttt{stop}} (stopped.north west);
\draw[tr] (paused_yes.north east) to[out=30,in=30] node[right, pos=0.38] {\scriptsize \texttt{stop}} (stopped.north east);
\end{tikzpicture}
}
\caption{State-transition diagram for dirty-tracking lifecycle states.}
\label{fig:dirty-tracking-state-machine}
\end{figure}

Transitions that violate owner checks, state preconditions, or snapshot guards
are rejected by the procfs handlers.

\subsubsection*{Permission Lifecycle}
Figure~\ref{fig:permission-lifecycle-flow} captures the permission policy used
to turn hardware write access into explicit software-recorded dirty events.

\paragraph{Start/Resume phase}
At \texttt{start} (and again at \texttt{resume}), the implementation walks the
owner VA space and revokes GPU write-class permissions from mapped pages,
forcing tracked mappings to \texttt{READ\_ONLY}. Before revocation, per-page
write/atomic masks are snapshot so \texttt{stop} can restore exact
pre-tracking permissions.

\paragraph{Active tracking phase}
During replayable fault servicing:
\begin{enumerate}
    \item Read faults remain mapped as \texttt{READ\_ONLY} and are not inserted
          into the dirty set.
    \item Write/atomic faults are serviced with writable permissions
          (\texttt{READ\_WRITE} or \texttt{READ\_WRITE\_ATOMIC}), and the page
          number plus timestamp are inserted into the live dirty table.
\end{enumerate}
This enforces first-write capture while preserving normal execution after fault
resolution.

\paragraph{Paused phase}
When paused, dirty recording is disabled, but fault handling still progresses:
read paths remain \texttt{READ\_ONLY}, while write/atomic paths may regain
writable permissions for correctness and forward progress. Because
\texttt{active=false}, no dirty inserts are emitted in this interval.

\paragraph{Cutover/Dump phase}
\texttt{query\_cutover} requires \texttt{snapshot\_pending==false}. Under the
query barrier, it swaps out the live table into a frozen snapshot and installs
a new live table so recording can continue. \texttt{query\_dump} then consumes
that frozen snapshot: in \texttt{delta} mode it emits snapshot contents
directly; in \texttt{cumulative} mode it merges snapshot entries into the
cumulative table before output.

\paragraph{Stop phase}
\texttt{stop} restores the previously captured per-page write/atomic permission
masks and tears down live, snapshot, and cumulative tracking state, returning
the system to its pre-tracking permission configuration.


\subsubsection*{Modified Driver Files}
All driver modifications reside under \texttt{kernel-open/nvidia-uvm/}.
The files containing functionality added or significantly modified by us are:

\begin{enumerate}
    \item \textbf{\texttt{uvm\_common.h} / \texttt{uvm\_common.c }}\\
          The primary implementation file.
          Defines \texttt{struct dirty\_page\_info} (page number, timestamp) and the three
          global dirty-set tables (\texttt{g\_dirty\_ds}, \texttt{g\_dirty\_ds\_snapshot},
          \texttt{g\_dirty\_ds\_cumulative}).
          Implements table lifecycle (\texttt{uvm\_dirty\_page\_table\_init},
          \texttt{uvm\_dirty\_page\_table\_destroy}), fault recording
          (\texttt{uvm\_dirty\_page\_table\_record}), snapshot management
          (\texttt{uvm\_dirty\_new\_snapshot}, \texttt{uvm\_dirty\_merge\_snapshot\_into\_cumulative\_locked}),
          and all eight \texttt{procfs} handler implementations including
          \texttt{uvm\_dirty\_procfs\_init}.
          Exports function pointers (\texttt{uvm\_dirty\_invalidate\_fn},
          \texttt{uvm\_dirty\_activate\_now\_fn}, \texttt{uvm\_dirty\_barrier\_end\_fn},
          \texttt{uvm\_dirty\_query\_barrier\_begin\_fn},
          \texttt{uvm\_dirty\_query\_barrier\_end\_fn}) that are registered from
          \texttt{uvm\_va\_space.c} at module init.\\

    \item \textbf{\texttt{uvm\_dirty\_ds.h}}\\
            Defines the \texttt{uvm\_dirty\_ds\_ops} interface
          (init, insert, lookup, for\_each\_in\_range, destroy), the \texttt{uvm\_dirty\_ds}
          struct holding ops pointer, backend-private data, and statistics, and inline timing
          wrappers for all operations.\\

    \item \textbf{\texttt{uvm\_dirty\_ds\_xarray.c}}\\
          xarray backend using atomic
          \texttt{xa\_cmpxchg} for first-write-wins insert.\\

    \item \textbf{\texttt{uvm\_dirty\_ds\_vector.c}}\\
          Sorted dynamic array backend with
          binary-search insert and lookup.\\

    \item \textbf{\texttt{uvm\_dirty\_ds\_vector\_unsorted.c}}\\
          Unsorted append-order vector
          with spinlock, preserving fault-arrival order.\\

    \item \textbf{\texttt{uvm\_dirty\_ds\_linked\_list.c}}\\
          Singly-linked list with sentinel node.\\

    \item \textbf{\texttt{uvm\_dirty\_ds\_bitmap.c}}\\
          Flat bit array covering $2^{38}$ pages
          (32\,MB fixed allocation); O(1) operations; timestamps not stored.

    \item \textbf{\texttt{uvm\_dirty\_ds\_nested\_bitmap.c}}\\
          Two-level hierarchical bitmap
          for sparse page sets.\\

    \item \textbf{\texttt{uvm\_dirty\_ds\_chunked.c}}\\
          Chunked vector with a background
          kthread sorter for deferred ordering.\\

    \item \textbf{\texttt{uvm\_gpu\_replayable\_faults.c}}\\
          The write-fault hook.
          After standard fault servicing, checks \texttt{uvm\_dirty\_tracking\_active()}
          and \texttt{service\_access\_type >= UVM\_FAULT\_ACCESS\_TYPE\_WRITE}, then
          calls \texttt{uvm\_dirty\_page\_table\_record} with the faulting page number
          and \texttt{ktime\_get\_ns()} timestamp.\\

    \item \textbf{\texttt{uvm\_va\_space.c}}\\
          Implements the barrier infrastructure.
          \texttt{uvm\_dirty\_downgrade\_all\_permissions} walks all va-ranges and
          va-blocks for the owner's va\_space, calling
          \texttt{uvm\_va\_block\_revoke\_prot\_mask} with \texttt{UVM\_PROT\_READ\_WRITE}
          to strip write access from all GPU PTEs, followed by \texttt{uvm\_tracker\_wait}
          to confirm GPU processing.
          \texttt{uvm\_dirty\_query\_barrier\_begin} and
          \texttt{uvm\_dirty\_query\_barrier\_end} implement the cutover barrier by
          locking all parent GPU replayable-fault ISRs and spinning until the fault
          buffer is empty.
          \texttt{uvm\_va\_space\_dirty\_init} registers all five function pointers
          at module init.\\

    \item \textbf{\texttt{uvm\_va\_space.h}}\\
          Declares \texttt{uvm\_va\_space\_dirty\_init},
          the registration hook invoked during module initialisation.\\

    \item \textbf{\texttt{uvm\_va\_block.c}}\\
          \texttt{compute\_new\_permission} modified to
          force \texttt{READ\_ONLY} on read faults when tracking is active for the
          faulting va\_space's owner, preventing read faults from obtaining write
          mappings that would bypass recording.\\

    \item \textbf{\texttt{uvm\_migrate.c}}\\
          \texttt{block\_migrate\_map\_unmapped\_pages}
          modified to map GPU-migrated pages as \texttt{READ\_ONLY} when tracking is
          active, closing the \texttt{cudaMemPrefetchAsync} bypass.\\

    \item \textbf{\texttt{uvm\_va\_block.h}}\\
          Extends per-GPU va-block state with
          \texttt{dirty\_downgraded\_pages} and
          \texttt{dirty\_downgraded\_atomic\_pages} masks used to restore exact
          pre-tracking permissions at stop.\\

    \item \textbf{\texttt{uvm\_procfs.c}}\\
          Integrates the dirty-tracking \texttt{procfs}
          interface into \texttt{/proc/driver/nvidia-uvm/} by invoking
          \texttt{uvm\_dirty\_procfs\_init(uvm\_proc\_dir)} during procfs setup.\\

    \item \textbf{\texttt{nvidia-uvm-sources.Kbuild}}\\
          Adds all dirty data-structure backend
          source files (\texttt{uvm\_dirty\_ds\_xarray.c},
          \texttt{uvm\_dirty\_ds\_bitmap.c},
          \texttt{uvm\_dirty\_ds\_vector.c},
          \texttt{uvm\_dirty\_ds\_vector\_unsorted.c},
          \texttt{uvm\_dirty\_ds\_chunked.c},
          \texttt{uvm\_dirty\_ds\_nested\_bitmap.c},
          \texttt{uvm\_dirty\_ds\_linked\_list.c}) to the module build list.\\

    \item \textbf{\texttt{uvm.c}}\\
          Top-level integration; calls
          \texttt{uvm\_va\_space\_dirty\_init} during module initialisation so
          callback pointers are registered before tracking control paths execute.
\end{enumerate}

\subsection*{Procfs Interface}

Eight procfs entries are created under \texttt{/proc/driver/nvidia-uvm/}:

\begin{enumerate}
    \item \textbf{\texttt{dirty\_tracking\_start} (write)} \\ 
          Input: \texttt{"<pid> <delta|cumulative>"}\\
          Initializes a fresh live table for the target PID, selects query mode
          (\texttt{delta} or \texttt{cumulative}), executes permission downgrade
          over the owner's managed GPU mappings, and finally enables recording.\\
          The request is rejected if tracking is already started, if the input
          format is invalid, or if initialization/callback setup fails.\\

    \item \textbf{\texttt{dirty\_tracking\_stop} (write)} \\ 
          Input: \texttt{"<pid>"}\\
          Restores previously downgraded GPU page permissions, destroys all
          tracking state (live/snapshot/cumulative), clears snapshot-pending
          state, and flushes accumulated DS statistics to
          \texttt{/tmp/uvm\_dirty\_ds\_stats.txt}.\\
          The request is rejected if the supplied PID is not the active owner.\\

    \item \textbf{\texttt{dirty\_tracking\_pause} (write)} \\ 
          Input: \texttt{"<pid>"}\\
          Establishes the query barrier, sets \texttt{active=false}, and then
          releases the barrier. This ensures pause is synchronized with
          replayable-fault ISR state.\\
          The request is rejected if the PID does not match the active owner or
          if tracking is not currently active.\\

    \item \textbf{\texttt{dirty\_tracking\_resume} (write)} \\ 
          Input: \texttt{"<pid>"}\\
          Re-executes the same downgrade/activate/barrier-end sequence used by
          \texttt{start}, re-establishing deterministic write-fault capture
          after a paused interval.\\

    \item \textbf{\texttt{dirty\_tracking\_query\_cutover} (write)} \\ 
          Input: \texttt{"<pid>"}\\
          Acquires the query barrier, atomically swaps
          \texttt{live $\rightarrow$ snapshot} under write lock, installs a
          fresh live table, updates epoch metadata, sets
          \texttt{g\_snapshot\_pending=true}, and releases the barrier.\\
          The request is rejected if the previous snapshot has not yet been
          consumed by \texttt{query\_dump}.\\

    \item \textbf{\texttt{dirty\_tracking\_query\_dump} (read)} \\ 
          No input payload (read endpoint). Requires
          \texttt{g\_snapshot\_pending=true}.\\
          In \texttt{delta} mode, emits only the current frozen snapshot.
          In \texttt{cumulative} mode, first merges snapshot contents into
          \texttt{g\_dirty\_ds\_cumulative} (first-write-wins), then emits the
          cumulative union.\\
          Output includes a header line followed by one line per page:
          \texttt{0x\textit{addr}\ \textit{timestamp\_ns}}.\\
          Snapshot consumption is overflow-safe: when \texttt{seq\_file}
          overflows, snapshot state is retained and replayed on the next read.\\

    \item \textbf{\texttt{dirty\_ds\_stats\_toggle} (write)} \\ 
          Input: \texttt{"enable"} or \texttt{"disable"}\\
          Enables or disables instrumentation for DS and callback timing
          counters. Writing \texttt{enable} also resets all counters/timers to
          zero to provide a clean measurement window.\\

    \item \textbf{\texttt{dirty\_ds\_stats} (read)} \\ 
          Read-only statistics snapshot for the current session.\\
          Reports operation counts and average latency for
          init/insert/lookup/for-each/destroy, average lock-wait components,
          and callback timing for invalidate, activate-now, and barrier-end.
\end{enumerate}

\subsection*{Design Choices}
\subsubsection*{Determinism Guarantees}


\subsubsection*{Data Structures}
      \paragraph{Global state}
      Three \texttt{uvm\_dirty\_ds} instances hold the tracking data:
      \begin{enumerate}
            \item \texttt{g\_dirty\_ds} - live table; fault handler inserts here.
            \item \texttt{g\_dirty\_ds\_snapshot} - frozen copy created by cutover;
                  exists only between a cutover and the subsequent dump.
            \item \texttt{g\_dirty\_ds\_cumulative} - session-union table used in
                  cumulative mode; merge-accumulates snapshot contents at each dump.
      \end{enumerate}

      \paragraph{Backend abstraction}
      All three tables share the \texttt{uvm\_dirty\_ds\_ops} interface
      (\texttt{uvm\_dirty\_ds.h}), which provides init, insert, lookup,
      for\_each\_in\_range, and destroy.
      Inline wrappers measure operation latency when statistics are enabled.
      The active backend is selected at compile time by setting the \texttt{.ops}
      pointer in the static \texttt{g\_dirty\_ds} initialiser.


\subsubsection*{Concurrency Design}
A read-write semaphore (\texttt{uvm\_dirty\_ds\_rwsem}) protects all three
tables.
Fault-handler insertions and lookups take the read side concurrently.
Lifecycle operations (init, destroy, cutover swap) take the write side
exclusively, blocking new insertions only during the brief pointer swap.
A separate mutex (\texttt{uvm\_dirty\_lifecycle\_lock}) serialises the procfs
lifecycle handlers (start, stop, pause, resume, cutover) against each other.
